%% Full order, Frobenius, equality, and inhabitance arguments.

\section{First-order predicate doctrines}

\begin{annexitem}{annex:p10.quant.1}{thm:p10-doctrine-quantifier-laws}{Product substitution squares and adjoint order laws}
\begin{proof}
Write \(\pi_A:I\times A\to A\) for the second product projection.
For \(f:J\to I\), suppose \(a:K\to J\) and
\(b:K\to I\times A\) satisfy \(f a=\pi_{I,A}b\).
The unique mediating map into \(J\times A\) is
\[
                              h=\langle a,\pi_A b\rangle.
\]
Its first projection is \(a\), and
\[
               (f\times1_A)h
                  =\langle f a,\pi_A b\rangle
                  =\langle\pi_{I,A}b,\pi_A b\rangle=b.
\]
Conversely, these two requirements prescribe both projections
of any mediator, hence prescribe \(h\) uniquely.
The parameter square is therefore a pullback.

For uniqueness of a left adjoint \(L\) to a fixed monotone
map \(R\), the values \(L(a)\) are characterized by
\[
                              L(a)\leq b\quad\Longleftrightarrow\quad
                              a\leq R(b).
\]
Two such values have the same upper bounds; taking each as
an upper bound of itself proves mutual inequalities.
Antisymmetry identifies them.
Right adjoints are similarly characterized by their lower
bounds and are unique.

For \(L\dashv R\), substituting \(b=L(a)\) in the adjunction
gives \(a\leq R L(a)\), and substituting \(a=R(b)\) gives
\(L R(b)\leq b\).
Apply these observations first to
\(\exists_\pi\dashv\pi^*\) and then to
\(\pi^*\dashv\forall_\pi\).
They are exactly the four displayed unit/counit inequalities.

For every \(\beta\in P(I)\), adjunction and the join
universal property give
\[
 \begin{aligned}
 \exists_\pi(\alpha\vee\alpha')\leq\beta
 &\ \Longleftrightarrow\
       \alpha\vee\alpha'\leq\pi^*\beta\\
 &\ \Longleftrightarrow\
       \alpha\leq\pi^*\beta\ \text{and}\ \alpha'\leq\pi^*\beta\\
 &\ \Longleftrightarrow\
       \exists_\pi\alpha\vee\exists_\pi\alpha'\leq\beta.
 \end{aligned}
\]
The two left-hand values have the same upper bounds, so
they are equal.  Moreover \(\bot\leq\pi^*\beta\) for every
\(\beta\), hence \(\exists_\pi\bot\) is below every
\(\beta\) and is bottom.

For meets, test lower bounds instead:
\[
 \begin{aligned}
 \beta\leq\forall_\pi(\alpha\wedge\alpha')
 &\ \Longleftrightarrow\
       \pi^*\beta\leq\alpha\wedge\alpha'\\
 &\ \Longleftrightarrow\
       \pi^*\beta\leq\alpha\ \text{and}\ \pi^*\beta\leq\alpha'\\
 &\ \Longleftrightarrow\
       \beta\leq\forall_\pi\alpha\wedge\forall_\pi\alpha'.
 \end{aligned}
\]
Since \(\pi^*\beta\leq\top\) for every \(\beta\),
\(\forall_\pi\top=\top\).

Finally suppose the fibres are Boolean.
Reindexing preserves negation, because
\(\neg\beta=(\beta\Rightarrow\bot)\) and it preserves
both operations.  Boolean complementation and adjunction give
\[
 \begin{aligned}
 \beta\leq\neg\exists_\pi\neg\alpha
 &\ \Longleftrightarrow\
       \exists_\pi\neg\alpha\leq\neg\beta\\
 &\ \Longleftrightarrow\
       \neg\alpha\leq\pi^*\neg\beta=\neg\pi^*\beta\\
 &\ \Longleftrightarrow\
       \pi^*\beta\leq\alpha\\
 &\ \Longleftrightarrow\
       \beta\leq\forall_\pi\alpha .
 \end{aligned}
\]
Equality of the two values follows by their identical
lower bounds.
None of the adjoint calculations invoked a Beck--Chevalley
identity.
\end{proof}
\end{annexitem}

\begin{annexitem}{annex:p10.quant.2}{thm:p10-doctrine-frobenius}{The Heyting implication argument}
\begin{proof}
Fix \(a\in H\), \(b\in K\), and test an arbitrary upper
bound \(c\in K\).
The adjunction, Heyting residuation in the two algebras,
and preservation of implication give
\[
 \begin{aligned}
 L(a\wedge R(b))\leq c
 &\ \Longleftrightarrow\ a\wedge R(b)\leq R(c)\\
 &\ \Longleftrightarrow\ a\leq(R(b)\Rightarrow R(c))\\
 &\ \Longleftrightarrow\ a\leq R(b\Rightarrow c)\\
 &\ \Longleftrightarrow\ L(a)\leq b\Rightarrow c\\
 &\ \Longleftrightarrow\ L(a)\wedge b\leq c.
 \end{aligned}
\]
Thus \(L(a\wedge R(b))\) and \(L(a)\wedge b\) have the
same upper bounds.  Using each value in turn as \(c\)
proves both inequalities, and antisymmetry gives equality.
For the doctrine, take
\[
 H=P(I\times A),\quad K=P(I),\quad
 L=\exists_\pi,\quad R=\pi^*.
\]
The defining reindexing homomorphism preserves implication,
so all the hypotheses hold.
The residuation and implication-preservation steps are the
specific reasons this argument does not follow from
finite-meet preservation alone.
\end{proof}
\end{annexitem}

\begin{annexitem}{annex:p10.quant.3}{thm:p10-doctrine-equality-laws}{Parameterized replacement and the equality laws}
\begin{proof}
\paragraph{Reflexivity with the actual terminal product.}
The unit of \(L_e\dashv e^*\), applied to
\(\top_{I\times A}\), gives
\[
                 \top_{I\times A}\leq e^*r^*\delta_A.
\]
Indeed \(L_e(\top)=r^*\delta_A\), since reindexing preserves
top and meets.
For \(I=1\), the composite \(r e:1\times A\to A\times A\)
is \(\Delta_A\pi_A\).
Let \(s_A=\langle!_A,1_A\rangle:A\to1\times A\).
Reindexing the displayed inequality along \(s_A\), and using
\(\pi_A s_A=1_A\), yields
\[
                              \top_A\leq\Delta_A^*\delta_A.
\]
The reverse inequality holds because \(\top_A\) is greatest.
This proves reflexivity without identifying \(1\times A\)
literally with \(A\).

\paragraph{Replacement in every predicate with parameters.}
The product equations give
\[
                         q_{I,A}e_{I,A}=1_{I\times A}.
\]
Hence
\(\gamma\leq e^*q^*\gamma=\gamma\).
The equality adjunction applied to this inequality gives
\[
                          p^*\gamma\wedge r^*\delta_A\leq q^*\gamma.
                                                                  \tag{1}
\]
This is the replacement law with the parameter object retained.

\paragraph{Symmetry.}
In (1) take \(I=A\) and
\(\gamma=\tau_A^*\delta_A\in P(A\times A)\).
Define
\[
 m=\langle\Delta_A\pi_1,\pi_2\rangle:
                           A\times A\longrightarrow(A\times A)\times A.
\]
Comparing product projections gives
\[
                         p m=\Delta_A\pi_1,\qquad
                         r m=1_{A\times A},\qquad
                         q m=1_{A\times A}.
\]
Also \(\tau_A\Delta_A=\Delta_A\).
After reindexing (1) along \(m\), its first conjunct is
\[
             (\Delta_A\pi_1)^*\tau_A^*\delta_A
                        =\pi_1^*\Delta_A^*\delta_A=\top_{A\times A}.
\]
The resulting inequality is
\(\delta_A\leq\tau_A^*\delta_A\).
Apply \(\tau_A^*\) and use \(\tau_A^2=1_{A\times A}\)
to obtain the opposite inequality.  Thus equality is symmetric.

\paragraph{Transitivity.}
In (1) take \(I=A\) and \(\gamma=\delta_A\).
The three pairs selected by \(p,r,q\) are respectively
the first/second, second/third, and first/third copies of \(A\).
The result is exactly the transitivity inequality in the statement.

\paragraph{Preservation by arrows.}
Let \(f:A\to B\), and put
\(\gamma=(f\times f)^*\delta_B\in P(A\times A)\).
The product equations imply
\((f\times f)\Delta_A=\Delta_B f\), so reflexivity in \(B\)
gives
\[
                         \Delta_A^*\gamma
                               =f^*\Delta_B^*\delta_B=\top_A.
\]
Take \(I=A\) in (1) and reindex along the same \(m\)
as in the symmetry proof.  The first conjunct becomes
\(\pi_1^*\Delta_A^*\gamma=\top\), and \(r m=q m=1\).
Thus \(\delta_A\leq\gamma=(f\times f)^*\delta_B\).

Reindexing these inequalities along tuples of term arrows
preserves their order, meets, and top.  It therefore gives
the asserted equality laws for interpreted terms and predicates.
No step requires a faithful interpretation of base arrows or
a representation of predicates by subobjects.
\end{proof}
\end{annexitem}

\begin{annexitem}{annex:p10.quant.4}{thm:p10-doctrine-inhabited}{Inhabitance and removal of an unused variable}
\begin{proof}
For \(!_I:I\to1\), existential Beck--Chevalley gives
\[
 \exists_{\pi_{I,A}}(!_I\times1_A)^*\top_{1\times A}
       =!_I^*\exists_{\pi_{1,A}}\top_{1\times A}.
\]
Reindexing preserves top.  If \(A\) is inhabited, the
right side is \(!_I^*\top_1=\top_I\), and the left side
is \(\exists_{\pi_{I,A}}\top_{I\times A}\).
This proves inhabitance in every parameter context.
The converse takes \(I=1\).

If \(a:1\to A\) is a global element, then
\(s=\langle1_1,a\rangle:1\to1\times A\) is a section of
\(\pi=\pi_{1,A}\).
The existential unit at top is
\[
                    \top_{1\times A}\leq
                         \pi^*\exists_\pi\top_{1\times A}.
\]
Apply \(s^*\), which preserves top, and use \(\pi s=1_1\).
It gives
\(\top_1\leq\exists_\pi\top_{1\times A}\).
Thus \(A\) is inhabited.

Fix any parameter object \(I\) and put \(\pi=\pi_{I,A}\).
Frobenius and the just-proved parameter form of inhabitance give
\[
 \exists_\pi\pi^*\beta
    =\exists_\pi(\top_{I\times A}\wedge\pi^*\beta)
    =(\exists_\pi\top_{I\times A})\wedge\beta
    =\beta.
                                                                  \tag{1}
\]
If \(\pi^*\alpha\leq\pi^*\beta\), apply the monotone
\(\exists_\pi\) and use (1) to obtain \(\alpha\leq\beta\).
The reverse implication is monotonicity of reindexing.
This proves order reflection.

The universal unit gives \(\beta\leq\forall_\pi\pi^*\beta\).
The counit at \(\pi^*\beta\) gives
\[
                  \pi^*\forall_\pi\pi^*\beta\leq\pi^*\beta.
\]
Order reflection gives
\(\forall_\pi\pi^*\beta\leq\beta\), proving the other equality.

For a semantic sequent, meet preservation makes extension by an
unused variable reindex both the conjunction of its assumptions
and its conclusion along \(\pi\).
The equivalence of the two inequalities is precisely the
proved order-reflection statement.
For a finite list of inhabited unused variables, apply it to
the successive context projections in reverse order.
The assumptions are preserved because inhabitance was proved
for every parameter context.
\end{proof}
\end{annexitem}
